\section{Review of Relevant Literature}
The modeling of financial market volatility began with the autoregressive conditional heteroskedasticity (ARCH) model introduced by \citet{engle1982autoregressive}, which was subsequently generalized to the GARCH model by \citet{bollerslev1986generalized}. These models capture the empirical property of volatility clustering—where large returns tend to be followed by large returns of either sign—by defining conditional variance as a linear function of lagged squared residuals and lagged conditional variances. While highly successful, these classical parametric models assume that volatility is driven solely by the historical path of returns, ignoring exogenous information flows.

To address this gap, researchers have integrated exogenous variables directly into the variance equation, leading to the GARCH-X class of models. In recent years, textual sentiment extracted from newspapers and social media has emerged as a primary exogenous input. Early works utilized simple bag-of-words approaches or dictionary-based methods to measure sentiment. However, these methods fail to capture context, negation, or financial domain nuances. 

The advent of transformer-based language models, such as BERT, revolutionized financial NLP. In the context of the Vietnamese language, \citet{vu2023domain} demonstrated that domain-adaptive masked language modeling (MLM) on financial texts significantly improves the performance of sentiment classifiers. When pre-trained transformer models are fine-tuned on Vietnamese financial sitemaps, they can effectively classify headlines into positive, negative, or neutral sentiment categories, capturing contextual relationships unique to corporate and macroeconomic disclosures.

However, incorporating sentiment scores directly into linear econometric models often yields mixed results. As argued by \citet{leber2025nonlinear}, the relationship between news sentiment and index volatility is highly non-linear and asymmetrical. Systemic responses to news shocks often depend on current market states and are subject to thresholds, where mild news has little effect but extreme positive or negative sentiment triggers disproportionate retail trading activity. Linear GARCH-X models struggle to adapt to these non-linearities and frequently suffer from parameter instability and out-of-sample forecast degradation.

Consequently, recent literature has explored hybrid architectures that combine parametric econometrics with non-linear machine learning models. In these hybrid models, the GARCH stage filters out the linear time-series dynamics of returns, while a recurrent neural network—such as a Long Short-Term Memory (LSTM) network—is trained on the resulting GARCH residuals. The LSTM is uniquely suited for this task because it can process sequential features (such as rolling windows of daily returns and sentiment shares) and capture non-linear, state-dependent relationships. Our paper contributes to this growing literature by applying a hybrid GARCH + LSTM framework to the Vietnamese index, explicitly testing the model's robustness and comparing it to exogenous linear GARCH-X specifications.
